\documentclass[UTF8,a4paper,11pt]{ctexart}

\usepackage{amsmath,amssymb,mathtools}
\usepackage{geometry}
\usepackage{fancyhdr}
\usepackage{enumitem}
\usepackage{hyperref}

\geometry{left=2.5cm,right=2.5cm,top=2.5cm,bottom=2.5cm}
\setlength{\parindent}{2em}
\linespread{1.2}
\setlength{\headheight}{13.6pt}

\pagestyle{fancy}
\fancyhf{}
\fancyfoot[C]{\thepage}
\fancyhead[L]{数值分析\quad 第七次理论作业\quad 学号：241098117，姓名：张城宗}
\fancyhead[R]{习题四}

\newcommand{\dd}{\,\mathrm{d}}
\newcommand{\QED}{\hfill$\square$}

\begin{document}

\begin{center}
  {\LARGE\bfseries 数值分析\quad 第七次理论作业}\\[6pt]
  {\large 习题四第 1、3、7、11、13、41 题}\\[4pt]
  {\normalsize 学号：241098117\quad 姓名：张城宗}
\end{center}

\bigskip

%=====================================================
\noindent\textbf{习题四\quad 第 1 题}\quad
求 $A_1,A_2,A_3$，使得计算
\[
  I(f)=\int_{-1}^{1}f(x)\dd x
\]
的求积公式
\[
  I_2(f)=A_1f(-1)+A_2f\!\left(-\frac13\right)+A_3f\!\left(\frac13\right)
\]
的代数精确度至少为 $2$。

\medskip
\noindent\textbf{解}\quad
要求代数精确度至少为 $2$，就让公式对 $1,x,x^2$ 都准确成立。于是
\[
\begin{cases}
A_1+A_2+A_3=2,\\[2mm]
-A_1-\dfrac13 A_2+\dfrac13 A_3=0,\\[2mm]
A_1+\dfrac19 A_2+\dfrac19 A_3=\dfrac23.
\end{cases}
\]
由第三式减去第一式的 $\dfrac19$ 倍，或直接解这个方程组，得
\[
  A_1=\frac12,\qquad A_2=0,\qquad A_3=\frac32.
\]
所以所求公式为
\[
  \boxed{
  I_2(f)=\frac12 f(-1)+\frac32 f\!\left(\frac13\right).}
\]
顺便检验一下，若取 $f(x)=x^3$，精确积分为 $0$，而求积值为
\[
  -\frac12+\frac32\cdot\frac1{27}=-\frac49\ne0.
\]
所以它的代数精确度恰好为 $2$。

\bigskip\bigskip

%=====================================================
\noindent\textbf{习题四\quad 第 3 题}\quad
已知计算积分
\[
  \int_{-1}^{1} f(x)\dd x
\]
的求积公式
\[
  I_2(f)=C\,[f(x_1)+f(x_2)+f(x_3)]
\]
的代数精确度是 $3$。试确定系数 $C$ 和节点 $x_1,x_2,x_3$。

\medskip
\noindent\textbf{解}\quad
分别令 $f(x)=1,x,x^2,x^3$，由代数精确度为 $3$ 可得
\[
  3C=2,\qquad
  C(x_1+x_2+x_3)=0,
\]
\[
  C(x_1^2+x_2^2+x_3^2)=\frac23,\qquad
  C(x_1^3+x_2^3+x_3^3)=0.
\]
第一式给出
\[
  C=\frac23.
\]
因此
\[
  x_1+x_2+x_3=0,\qquad
  x_1^2+x_2^2+x_3^2=1,\qquad
  x_1^3+x_2^3+x_3^3=0.
\]
记三个节点的基本对称和为
\[
  s_1=x_1+x_2+x_3,\quad
  s_2=x_1x_2+x_1x_3+x_2x_3,\quad
  s_3=x_1x_2x_3.
\]
由 $s_1=0$ 及 $x_1^2+x_2^2+x_3^2=s_1^2-2s_2=1$，得
\[
  s_2=-\frac12.
\]
又
\[
  x_1^3+x_2^3+x_3^3=s_1^3-3s_1s_2+3s_3=3s_3,
\]
所以 $s_3=0$。于是节点应为方程
\[
  t^3+s_2t-s_3=0
\]
的三个根，即
\[
  t^3-\frac12t=0.
\]
故
\[
  \{x_1,x_2,x_3\}=\left\{-\frac1{\sqrt2},\,0,\,\frac1{\sqrt2}\right\}.
\]
节点顺序任意。最后得到
\[
  \boxed{
  C=\frac23,\qquad
  x_1=-\frac1{\sqrt2},\ x_2=0,\ x_3=\frac1{\sqrt2}.}
\]

\bigskip\bigskip

%=====================================================
\noindent\textbf{习题四\quad 第 7 题}\quad
对区间 $[a,b]$ 作等距剖分，节点为
\[
  x_0=a,\quad x_1,\ldots,x_n=b,\qquad
  h=\frac{b-a}{n},\quad x_i=a+ih.
\]
证明当 $n$ 为偶数时，
\[
  \int_a^b \omega_{n+1}(x)\dd x=0,\qquad
  \omega_{n+1}(x)=\prod_{i=0}^{n}(x-x_i),
\]
并说明此时闭 Newton-Cotes 型求积公式的代数精确度为 $n+1$。

\medskip
\noindent\textbf{证明}\quad
设 $n=2m$，并记区间中点为
\[
  c=\frac{a+b}{2}.
\]
因为是等距剖分，所以节点可以写成
\[
  x_{m+j}=c+jh,\qquad j=-m,-m+1,\ldots,m.
\]
令 $t=x-c$，则
\[
\begin{aligned}
  \omega_{n+1}(x)
  &=\prod_{j=-m}^{m}(t-jh)  \\
  &=t\prod_{j=1}^{m}(t^2-j^2h^2).
\end{aligned}
\]
这是关于 $t$ 的奇函数，而积分区间 $[a,b]$ 关于 $c$ 对称，即 $t\in[-mh,mh]$。因此
\[
  \int_a^b\omega_{n+1}(x)\dd x
  =\int_{-mh}^{mh}t\prod_{j=1}^{m}(t^2-j^2h^2)\dd t=0.
\]

下面看代数精确度。闭 Newton-Cotes 公式是由这 $n+1$ 个节点的 Lagrange 插值多项式积分得到的，因此它至少对所有不超过 $n$ 次的多项式精确。

若 $f(x)$ 是任意 $n+1$ 次多项式，设 $p_n(x)$ 为它在节点 $x_0,\ldots,x_n$ 上的插值多项式。由于 $f-p_n$ 在全部节点上为零，且次数不超过 $n+1$，故存在常数 $K$，使得
\[
  f(x)-p_n(x)=K\omega_{n+1}(x).
\]
由刚才证得的结论，
\[
  \int_a^b f(x)\dd x
  =\int_a^b p_n(x)\dd x+K\int_a^b\omega_{n+1}(x)\dd x
  =\int_a^b p_n(x)\dd x.
\]
而求积公式本来就是对 $p_n$ 的积分，所以对 $f$ 也准确成立。这说明当 $n$ 为偶数时，闭 Newton-Cotes 公式至少具有 $n+1$ 次代数精确度。

再看下一阶，一般取 $f(x)=(x-c)\omega_{n+1}(x)$ 时，所有节点处的函数值为零，求积值为 $0$；而
\[
  \int_a^b (x-c)\omega_{n+1}(x)\dd x
  =h^{n+3}\int_{-m}^{m}s^2\prod_{j=1}^{m}(s^2-j^2)\dd s
\]
不是闭 Newton-Cotes 公式的零误差矩。因此下一阶一般不再精确。于是闭 Newton-Cotes 型求积公式在 $n$ 为偶数时的代数精确度为
\[
  \boxed{n+1.}
\]
\QED

\bigskip\bigskip

%=====================================================
\noindent\textbf{习题四\quad 第 11 题}\quad
设函数 $f(x)$ 在 $[-h,h]$ 上充分可导，推导求积公式
\[
  \int_0^h f(x)\dd x\approx \frac h2[3f(0)-f(-h)]
\]
的余项。

\medskip
\noindent\textbf{解}\quad
记
\[
  R(f)=\int_0^h f(x)\dd x-\frac h2[3f(0)-f(-h)].
\]
这个公式对 $f(x)=1$ 和 $f(x)=x$ 都是准确的，所以余项从二阶导数开始。

用 Peano 核来写比较清楚。因为 $R$ 对一次多项式为零，
\[
  R(f)=\int_{-h}^{h}K(t)f''(t)\dd t,
\]
其中
\[
  K(t)=R\bigl((x-t)_+\bigr).
\]
分别计算：

当 $-h\le t<0$ 时，
\[
  K(t)=\int_0^h(x-t)\dd x-\frac h2[3(0-t)-0]
      =\frac{h(h+t)}2.
\]
当 $0\le t\le h$ 时，
\[
  K(t)=\int_t^h(x-t)\dd x
      =\frac{(h-t)^2}{2}.
\]
显然 $K(t)\ge0$。于是
\[
\begin{aligned}
  \int_{-h}^{h}K(t)\dd t
  &=\int_{-h}^{0}\frac{h(h+t)}2\dd t
    +\int_0^h\frac{(h-t)^2}{2}\dd t  \\
  &=\frac{h^3}{4}+\frac{h^3}{6}
   =\frac{5h^3}{12}.
\end{aligned}
\]
由积分中值定理，存在 $\xi\in(-h,h)$，使得
\[
  \boxed{
  R(f)=\frac{5h^3}{12}f''(\xi).}
\]
也就是说
\[
  \int_0^h f(x)\dd x
  =\frac h2[3f(0)-f(-h)]+\frac{5h^3}{12}f''(\xi).
\]

\bigskip\bigskip

%=====================================================
\noindent\textbf{习题四\quad 第 13 题}\quad
设基于点 $x_1,x_2,\ldots,x_{n+1}$ 的插值型求积公式
\[
  \sum_{i=1}^{n+1}A_if(x_i)
\]
具有 $m(m\ge n)$ 次代数精确度。证明对任意 $f(x)\in C[a,b]$，有
\[
  \sum_{i=1}^{n+1}A_if(x_i)=\int_a^b p_m(x)\dd x,
\]
其中 $p_m(x)$ 是 $f(x)$ 关于节点 $x_1,x_2,\ldots,x_{n+1}$ 的 $m$ 次插值多项式。

\medskip
\noindent\textbf{证明}\quad
由于 $p_m(x)$ 是关于这些节点的插值多项式，所以
\[
  p_m(x_i)=f(x_i),\qquad i=1,2,\ldots,n+1.
\]
于是求积公式作用在 $f$ 和 $p_m$ 上时，得到的离散值完全相同：
\[
  \sum_{i=1}^{n+1}A_if(x_i)
  =\sum_{i=1}^{n+1}A_ip_m(x_i).
\]
又因为求积公式具有 $m$ 次代数精确度，而 $p_m$ 的次数不超过 $m$，所以
\[
  \sum_{i=1}^{n+1}A_ip_m(x_i)=\int_a^b p_m(x)\dd x.
\]
两式合起来便得到
\[
  \boxed{
  \sum_{i=1}^{n+1}A_if(x_i)=\int_a^b p_m(x)\dd x.}
\]
当 $m>n$ 时，其余 $m-n$ 个插值条件可以取导数值相等等条件；上面的证明只用到了节点处函数值相等，以及公式对 $m$ 次多项式精确这两点。
\QED

\bigskip\bigskip

%=====================================================
\noindent\textbf{习题四\quad 第 41 题}\quad
确定数值微分公式
\[
  f'(x_0)\approx
  \frac{1}{2h}\left(3f(x_0+h)-f(x_0)-2f\!\left(x_0+\frac h2\right)\right)
\]
的截断误差表达式。

\medskip
\noindent\textbf{解}\quad
在 $x_0$ 处作 Taylor 展开：
\[
\begin{aligned}
  f(x_0+h)
  &=f(x_0)+hf'(x_0)+\frac{h^2}{2}f''(x_0)
    +\frac{h^3}{6}f^{(3)}(x_0)
    +\frac{h^4}{24}f^{(4)}(x_0)+O(h^5),\\
  f\!\left(x_0+\frac h2\right)
  &=f(x_0)+\frac h2f'(x_0)+\frac{h^2}{8}f''(x_0)
    +\frac{h^3}{48}f^{(3)}(x_0)
    +\frac{h^4}{384}f^{(4)}(x_0)+O(h^5).
\end{aligned}
\]
代入右端，记该差商为 $D_h$，则
\[
\begin{aligned}
  D_h
  &=\frac{1}{2h}\left(3f(x_0+h)-f(x_0)-2f\!\left(x_0+\frac h2\right)\right)\\
  &=f'(x_0)+\frac58h f''(x_0)
    +\frac{11}{48}h^2 f^{(3)}(x_0)
    +\frac{23}{384}h^3 f^{(4)}(x_0)+O(h^4).
\end{aligned}
\]
因此若把截断误差记为“近似值减真值”，则
\[
  \boxed{
  D_h-f'(x_0)
  =\frac58h f''(x_0)+\frac{11}{48}h^2 f^{(3)}(x_0)+O(h^3).}
\]
特别地，主项为
\[
  D_h-f'(x_0)=\frac58h f''(x_0)+O(h^2),
\]
所以该数值微分公式是一阶精度的。等价地，也可写成
\[
  f'(x_0)
  =
  \frac{1}{2h}\left(3f(x_0+h)-f(x_0)-2f\!\left(x_0+\frac h2\right)\right)
  -\frac58h f''(\xi),
  \quad \xi\in(x_0,x_0+h).
\]

\end{document}
